\section{Emergent Dirac and Lorentz}

The discrete walk reproduces relativistic propagation. The Dirac dispersion and a finite light cone are read off
a simulation of the rule rather than assumed. What follows is a measurement of how the substrate already moves,
not an importation of relativity into it.

A few terms are worth stating plainly. The \emph{Dirac equation} is the law for how an electron-like particle
moves when both quantum mechanics and special relativity apply at once, and its dispersion is the relation
between a wave's energy and its wavelength. \emph{Lorentz invariance} is the content of special relativity, that
the laws look the same to everyone moving at constant speed and that there is one finite top speed, light. A
\emph{light cone} is the reach of that top speed, the set of places a signal can get to from a given point.

\subsection{The light cone}

A pulse spreads one cell per beat. That is a finite, frame-independent maximum speed, $z = 1$, and it is the
light cone. For the massless case the spreading is exactly the discrete shift, a permutation of cell contents,
so continuity is not invoked anywhere in the statement. The light cone built from the $24$ directions is also
round, that is, isotropic at scale, a point taken up where isotropy is measured.

\begin{result}[A discrete light cone]
A pulse on the substrate spreads exactly one cell per beat \cite{pollard2026vibetest}. The maximum signal speed
is finite and frame-independent, with dynamical exponent $z = 1$. In the massless case the propagation is the
discrete shift itself, so the light cone is exact and uses no continuum limit.
\end{result}

\subsection{The Dirac dispersion}

Simulate the discrete two-component walk, a left mover and a right mover, and measure its dispersion $E(k)$ by
Fourier transforming the time series of a tracked cell. Two cases arise. In the massless case the rule is the
pure discrete shift, the measured dispersion is $E = k$ exactly, and the light cone is the shift itself with no
continuity assumed. In the massive case a coin rotation by the mass angle $m$ is added, the measured dispersion
obeys $\cos E = \cos m \cos k$, the lattice Dirac dispersion, and the rest mass is $E(0) = m$. Mass is the
coarse-grained mixing rate, nothing more.

\begin{result}[The lattice Dirac dispersion]
The discrete two-component walk has measured dispersion $E = k$ in the massless case and
$\cos E = \cos m \cos k$ in the massive case \cite{pollard2026vibetest}, the latter being the lattice Dirac
dispersion with rest mass $m$. The Dirac dispersion is thus the rule's own measured dispersion.
\end{result}

The reading is that the complex quantum amplitudes are the emergent description of the real discrete walk's
oscillation modes. The quantum layer is emergent. The base is the real discrete walk, and the amplitudes are how
that walk's modes are most economically written down. This is consistent with the cellular-automaton
reconstructions of the free Dirac field by D'Ariano and Bisio \cite{dariano2014, bisio2015}, where the same
lattice dispersion arises from a local, reversible update.

\subsection{Why masslessness needs an even-faced cell}

There is a clean geometric precondition for an exactly massless mode. The cell must have opposite faces. On an
odd-faced cell the instruction to go straight lands between two faces, which forces a small zigzag of order
$\pi/F$ at every step, an irreducible scatter. The lightest particle then has a floor mass of order
$(\pi/2F)^2$, and nothing is exactly massless. An even-faced cell with opposite pairs lets a signal go dead
straight, with mixing $\lambda = 0$, and so it hosts exactly massless modes. The $24$-cell is even-faced. Its
$24$ facets form $12$ opposite pairs, so $\{3,4,3,4\}$ permits the exactly-massless $E = k$ mode above, where a
generic odd tiling could not. The geometry is not cosmetic. The masslessness rests on the cell's antipodal
structure.

\begin{proposition}[The mass floor]
An exactly massless mode requires a cell with opposite faces. On an odd-faced cell with $F$ faces, going straight
forces a zigzag of order $\pi/F$ per step, giving every mode a floor mass of order $(\pi/2F)^2$. The $24$-cell's
$24$ facets form $12$ opposite pairs, so $\{3,4,3,4\}$ admits the exactly-massless $E = k$ mode.
\end{proposition}

The telegraph picture ties the cases together with a single knob. A persistent walk that scatters at rate
$\lambda$ is the telegrapher's equation, equivalently the Klein-Gordon or one-dimensional Dirac equation, with
mass proportional to $\lambda$. Setting $\lambda = 0$ gives the massless wave with $z = 1$. A finite $\lambda$
gives the relativistic massive case. Taking $\lambda$ to infinity recovers ordinary diffusion with $z = 2$. One
mechanism covers three limits.


This closes the bridge from the base to quantum mechanics. Relativistic quantum mechanics is an emergent limit
of the discrete substrate, with the massless light cone being exactly discrete, needing no continuum, and mass
emerging as a mixing rate.

\subsection{Caveat}

The cleanest reading is the one-dimensional, two-component demonstrator. The full $24$-direction version is the
bulk wave, and Lorentz invariance on the physical cusp carries a small cubic anisotropy that washes out in the
infrared. The precise statement is that the dispersion and the light cone are exact in the demonstrator and
isotropic in the limit, with the residual lattice anisotropy treated where isotropy and falsifiability are
discussed.
